% Generated draft: requires mathematical and source review.
\documentclass[11pt]{article}
\usepackage[margin=25mm]{geometry}
\usepackage{fontspec}
\setmainfont{Latin Modern Roman}
\usepackage{amsmath,amssymb,mathtools,unicode-math}
\setmathfont{Latin Modern Math}
\usepackage{enumitem,needspace,xurl,hyperref}
\setlength{\parindent}{0pt}
\setlength{\parskip}{5pt}
\setlength{\emergencystretch}{3em}
\newcommand{\N}{\mathbb{N}}
\newcommand{\Z}{\mathbb{Z}}
\newcommand{\Q}{\mathbb{Q}}
\newcommand{\R}{\mathbb{R}}
\newcommand{\C}{\mathbb{C}}
\newcommand{\F}{\mathbb{F}}
\newcommand{\A}{\mathbb{A}}
\newcommand{\PP}{\mathbb{P}}
\newcommand{\E}{\mathbb{E}}
\newcommand{\eps}{\varepsilon}
\newcommand{\dd}{\,\mathrm{d}}
\begin{document}
\section*{UnsolvedMath ID: 30003559}
\textbf{Source problem code:} catalogue entry
\subsection{Definitions and mathematical statement}
For a specified root datum with weight lattice \(P\) and simple coroots \(\{\alpha_i^\vee\}_{i\in I}\), its dominant weights are
\[
P_+=\{\lambda\in P:\langle\lambda,\alpha_i^\vee\rangle\geq 0\text{ for every }i\in I\}.
\]
This definition is conditional: the supplied record does not identify the root datum or the representation-theoretic setting. If that setting provides a finite-dimensional representation \(M_\lambda\) with weight decomposition \(M_\lambda=\bigoplus_{\mu\in P}(M_\lambda)_\mu\), its formal character and dimension are
\[
\operatorname{ch}M_\lambda=\sum_{\mu\in P}\dim (M_\lambda)_\mu\,e^\mu\in\Z[P],
\qquad
\dim M_\lambda=\sum_{\mu\in P}\dim (M_\lambda)_\mu,
\]
where \(\Z[P]\) is the group ring of \(P\), with formal basis elements \(e^\mu\). The record does not specify which representations \(M_\lambda\) are intended.

The supplied problem asks: ``Can one find Verlinde-like dimension formulas for other families of dominant weights, and do the resulting character formulas coincide with those predicted by Lusztig’s conjecture?'' In a fully specified setting this would require identifying families \(\mathcal F\subseteq P_+\), giving dimension formulas for the intended representations indexed by \(\lambda\in\mathcal F\), and determining whether their characters agree with the applicable Lusztig predictions throughout the stated range of parameters.

The phrase ``Verlinde-like'' is not defined, the baseline for ``other families'' is not identified precisely, and neither the applicable version of Lusztig's conjecture nor its hypotheses are supplied. Dimension formulas alone do not specify formal characters, so the construction of the ``resulting character formulas'' also requires clarification. Consequently, the supplied evidence does not determine a complete quantified mathematical conjecture.

The supplied dated curation says that the dimension question is answered for the ``p minus 1 fundamental weights,'' while the broader family question remains open. It does not define \(p\) or identify these weights. It also says that the related general Lusztig conjecture is false but that this does not settle the specific character comparison. These are supplied qualifications, not independently verified mathematical conclusions.
\subsection{Short English statement}
The problem asks for Verlinde-like dimension formulas for further families of dominant weights and for a comparison of associated character formulas with Lusztig's predictions. The group or root datum, representations, parameter restrictions, meaning of ``Verlinde-like,'' and intended families are missing, so a precise conjecture cannot yet be extracted. The supplied status is partially solved; this statement-only draft requires human mathematical and source review.
\subsection{Sources}
\begin{itemize}
\item Local UnsolvedMath registry: permanent ID \(30003559\), source problem code \texttt{catalogue entry}, title ``Verlinde Formulas and Lusztig Characters for New Weight Families,'' in the Oberwolfach Reports Open Problems collection. The record gives proposed year \(2017\) and names no proposer. It labels the statement \texttt{exact} and describes it as self-contained and publication-ready; these are supplied registry assertions, not independent verification.
\item Supplied report citation: \emph{Low-dimensional Topology and Number Theory} (2018), Oberwolfach Reports, pp.\ 2363--2426, DOI \url{https://doi.org/10.4171/owr/2017/38}. The background attributes the problem to workshop open-problem discussions documented by Oberwolfach Reports. No problem-specific page or passage is supplied.
\item Supplied literature source, labeled ``Low dimensional Topology and Number Theory'': \url{https://doi.org/10.4171/owr/2017/38}.
\item Supplied literature source, labeled ``Light leaves and Lusztig conjecture'': \url{https://doi.org/10.1016/j.aim.2015.04.022}.
\item Literature assessment dated exactly \texttt{2026-08-21}: the supplied curation labels the problem partially solved, reports dimension formulas for the ``p minus 1 fundamental weights,'' leaves the broader family question open, and says that failure of the related general Lusztig conjecture does not settle this particular character comparison. This assessment is supplied dated curation derived from the OWR material, not independently verified. The record explicitly requests independent verification of time-sensitive status and cited claims.
\end{itemize}
\end{document}
